% Part: computability
% Chapter: recursive-functions
% Section: composition

\documentclass[../../../include/open-logic-section]{subfiles}

\begin{document}

\olfileid{cmp}{rec}{com}
\olsection{संयोजन}

$f$ आणि $g$ ही नैसर्गिक संख्यांची दोन एकस्थानी फलने असतील, तर त्यांचे
संयोजन करता येते: $h(x) = f(g(x))$. मग नवीन फलन~$h(x)$ हे $f$ आणि~$g$
या फलनांपासून \emph{संयोजनाने} परिभाषित होते. याचे एकाहून अधिक आदाने
असलेल्या फलनांसाठी सामान्यीकरण करायचे आहे.

ते करण्याचा एक मार्ग असा: $f$ हे $k$-स्थानी फलन आहे आणि $g_0$,
\dots, $g_{k-1}$ ही सर्व $n$-स्थानी असलेली $k$ फलने आहेत असे समजा.
मग पुढीलप्रमाणे नवीन $n$-स्थानी फलन $h$ परिभाषित करता येते:
\[
h(x_0, \dots, x_{n-1}) =
f(g_0(x_0, \dots, x_{n-1}), \dots, g_{k-1}(x_0, \dots, x_{n-1}))
\]
$f$ आणि सर्व $g_i$ संगणनक्षम असतील, तर $h$ देखील संगणनक्षम असते:
$h(x_0, \dots, x_{n-1})$ संगणित करण्यासाठी प्रथम प्रत्येक $i = 0$,
\dots,~$k-1$ साठी $y_i = g_i(x_0, \dots, x_{n-1})$ ही मूल्ये संगणित
करा. नंतर ही मूल्ये $f$ ला आदाने म्हणून देऊन
$h(x_0, \dots, x_{n-1}) = f(y_0, \dots, y_{k-1})$ संगणित करा.

आधीपासून उपलब्ध असलेल्या काही फलनांचा उपयोग करून नवीन फलन संगणित करताना
जे घडते त्याचे हे लक्षणचित्रण अनावश्यकपणे बंधनकारक वाटू शकते. एक कारण
म्हणजे, कधीकधी आपण फलनाची सर्व आदाने वापरत नाही; उदाहरणार्थ,~$\Add$
च्या आदिम पुनरावर्ती व्याख्येत वापरण्यासाठी $g(x, y, z) = \Succ(z)$
परिभाषित केले तेव्हा तसेच झाले. पुढील फलने वापरण्याची आपल्याला मुभा आहे
असे समजा:
\[
\Proj{n}{i}(x_0, \dots, x_{n-1}) = x_i
\]
फलनांना~$\Proj{k}{i}$ \emph{प्रक्षेपण} फलने म्हणतात:
$\Proj{n}{i}$ हे $n$-स्थानी फलन आहे. मग $g$ ची व्याख्या अशी देता येते:
\[
g(x, y, z) = \Succ(\Proj{3}{2}(x, y, z)).
\]
येथे $f$ ची भूमिका $1$-स्थानी फलन $\Succ$ बजावते, म्हणून $k=1$.
तसेच $g_0$ ची भूमिका बजावणारे एक $3$-स्थानी फलन $\Proj{3}{2}$ आहे.
परिणाम म्हणजे तिसऱ्या आदानाचा उत्तरवर्ती परत करणारे $3$-स्थानी फलन.

प्रक्षेपण फलनांमुळे आदानांचा क्रम बदलून किंवा काही आदाने एकरूप करूनही
नवीन फलने परिभाषित करता येतात. उदाहरणार्थ, फलन $h(x) = \Add(x, x)$ ची
व्याख्या अशी देता येते:
\[
h(x_0) = \Add(\Proj{1}{0}(x_0),\Proj{1}{0}(x_0)).
\]
येथे $k=2$, $n=1$, $f(y_0,y_1)$ ची भूमिका $\Add$ बजावते, आणि
$g_0(x_0)$ व $g_1(x_0)$ या दोहोंची भूमिका~$\Proj{1}{0}(x_0)$ हे
एकस्थानी प्रक्षेपण फलन (म्हणजे एकरूपता फलन) बजावते.

$f(y_0, y_1)$ हे फलन आपल्याकडे आधीच असेल, तर फलन
$h(x_0, x_1) = f(x_1, x_0)$ ची व्याख्या अशी देता येते:
\[
h(x_0, x_1) = f(\Proj{2}{1}(x_0, x_1),\Proj{2}{0}(x_0, x_1)).
\]
येथे $k=2$, $n = 2$, आणि $g_0$ व $g_1$ यांच्या भूमिका अनुक्रमे
$\Proj{2}{1}$ आणि~$\Proj{2}{0}$ बजावतात.

$g_0$, \dots,~$g_{k-1}$ या सर्वांची स्थानसंख्या~$n$ समान असणे आवश्यक
आहे, याबद्दलही शंका वाटू शकते. (लक्षात ठेवा, फलनाची \emph{स्थानसंख्या}
म्हणजे त्याच्या आदानांची संख्या; $n$-स्थानी फलनाची स्थानसंख्या~$n$
असते.) पण प्रक्षेपण फलने समाविष्ट केल्याने अपेक्षित लवचीकता मिळते.
उदाहरणार्थ, $f$ आणि~$g$ ही $3$-स्थानी फलने आहेत आणि पुढीलप्रमाणे
परिभाषित केलेले $h$ हे $2$-स्थानी फलन आहे असे समजा:
\[
h(x,y) = f(x,g(x,x,y),y).
\]
प्रक्षेपण फलने वापरून~$h$ ची व्याख्या पुढीलप्रमाणे पुन्हा लिहिता येते:
\[
h(x,y) = f(\Proj{2}{0}(x,y), g(\Proj{2}{0}(x,y), \Proj{2}{0}(x,y),
\Proj{2}{1}(x,y)), \Proj{2}{1}(x,y)).
\]
मग $h$ हे $f$ चे $\Proj{2}{0}$, $l$ आणि $\Proj{2}{1}$ यांच्याशी
केलेले संयोजन आहे, येथे
\[
l(x,y) = g(\Proj{2}{0}(x,y),\Proj{2}{0}(x,y),\Proj{2}{1}(x,y)),
\]
म्हणजे $l$ हे $g$ चे $\Proj{2}{0}$, $\Proj{2}{0}$ आणि~$\Proj{2}{1}$
यांच्याशी केलेले संयोजन आहे.

\end{document}
